\section{What finite arithmetic constraints do not determine}\label{sec:controls}

The estimates above use the physical theta coefficients, all equal to one.
It is useful to identify which of their premises would survive a small
change of coefficients. We prove a collision construction within a precise
class of positive arithmetic sources. The construction preserves any
specified finite set of prime identities and linear measurements, but
fails the full modular reflection. The distinction can be proved directly.

\subsection{Rigidity of the complete reflection law}

We first allow even real coefficients $(a_n)_{n\in\Z}$ of at most polynomial
growth and set
\[
 \Theta_a(y)=\sum_{n\in\Z}a_ne^{-\pi n^2y},\quad
 \psi_a(u)=\sum_{n\ge1}a_np_n(u),\quad
 \Phi_a(u)=\sum_{n\ge1}a_n\phi_n(u).
\]
Polynomial growth leaves all the local convergence arguments of
Section~\ref{sec:theta} valid. Define
$P_a(u)=a_0e^u+2\psi_a(u)$. Then $P_a''-P_a=16\Phi_a$.

\begin{theorem}[Coefficient rigidity]\label{thm:rigidity}
If
\begin{equation}\label{eq:weighted-reflection}
 \Theta_a(y)=y^{-1/2}\Theta_a(1/y)\qquad(y>0),
\end{equation}
then $a_n=a_0$ for every $n\in\Z$.
\end{theorem}
\begin{proof}
We give the distribution argument, including its uniqueness step.
A tempered distribution is a continuous linear functional on the Schwartz
space $\Sch(\R)$. Polynomial growth makes
$\mu=\sum_{n\in\Z}a_n\delta_n$ such a distribution. Its Fourier transform is
defined by $\langle\widehat\mu,f\rangle=\langle\mu,\widehat f\rangle$,
with the convention of Section~\ref{sec:theta}. Put $\nu=\mu-\widehat\mu$.
The Gaussian transform formula and \eqref{eq:weighted-reflection} give
\[
 \langle\nu,e^{-\pi yv^2}\rangle
 =\Theta_a(y)-y^{-1/2}\Theta_a(1/y)=0 .
\]
Both $\mu$ and $\widehat\mu$ are even. Differentiating this identity at
$y=1$ in the Schwartz topology shows that
$\langle\nu,v^{2k}e^{-\pi v^2}\rangle=0$ for all $k\ge0$; the odd
moments vanish by evenness.

The function
$J(z)=\langle\nu,e^{-\pi v^2+zv}\rangle$ is entire. To justify this,
on every compact $z$ set each Schwartz seminorm of the integrand and
its $z$ derivatives is bounded by a polynomial in $|v|$ times
$e^{-\pi v^2+C|v|}$. Differentiation is therefore valid in that topology.
All Taylor coefficients of $J$ at zero vanish, so $J=0$.
For $z=-2\pi i\xi$, this says that the Fourier transform of the
distribution $e^{-\pi v^2}\nu$ vanishes. More explicitly, integrate
$J(-2\pi i\xi)$ against an arbitrary Schwartz test function and interchange
the integral and distribution pairing; the same seminorm bounds justify
the interchange. Fourier inversion on Schwartz functions, and hence
injectivity on tempered distributions, imply $e^{-\pi v^2}\nu=0$.
For a compactly supported smooth test function $f$, division by the
nonvanishing Gaussian preserves smooth compact support. Consequently
$\langle\nu,f\rangle=0$. Such tests determine a distribution, so $\nu=0$.

We have proved $\mu=\widehat\mu$. On the other hand,
$\widehat\mu=\sum_na_ne^{-2\pi inv}$ is a $1$-periodic distribution:
translation by one fixes each term and the distributional sum. Thus
$\mu$ is $1$-periodic. Testing in small neighborhoods of two consecutive
integers gives $a_{n+1}=a_n$, which proves the assertion.
\end{proof}

The primitive endpoint supplies exactly the normalization lost by applying
$\partial_u^2-1$. Indeed, put $D(u)=P_a(u)-P_a(-u)$. Then
\[
 D''-D=16(\Phi_a(u)-\Phi_a(-u)),\quad D(0)=0,\quad D'(0)=2P_a'(0).
\]
Solving this elementary initial value problem gives
\begin{equation}\label{eq:defect}
 D(u)=2P_a'(0)\sinh u+
 16\int_0^u\sinh(u-v)(\Phi_a(v)-\Phi_a(-v))\dd v .
\end{equation}
This formula follows either by differentiating the right side twice or
by variation of constants. If $\Phi_a$ is even and
$\psi_a'(0)=-a_0/2$, then $P_a'(0)=0$ and \eqref{eq:defect} gives $D=0$.
The change of variable $y=e^{4u}$ then gives
\eqref{eq:weighted-reflection}, so Theorem~\ref{thm:rigidity} applies.
In particular, for a signed perturbation with $a_0=0$, source evenness
and $\psi_a'(0)=0$ force every coefficient to vanish.

\subsection{Prime orbits and finite measurements}

Fix a finite prime set $\mathcal P$, an integer $N\ge1$, and finitely many
real linear measurements of a source and its primitive. Each measurement
is assumed to be defined on the span of the theta source and the orbit
sources below. Examples are finitely
many source derivatives at specified real points, transform derivatives
at fixed time and spectral points, and convergent weighted integrals.
A complex measurement means its real and imaginary parts separately.
No continuity assumption is needed for the finite-dimensional construction;
continuity matters only if one subsequently passes to a limit.

Let
\[
 \mathcal S(\mathcal P)=\left\{\prod_{p\in\mathcal P}p^{k_p}:
                              k_p\in\Z_{\ge0}\right\}.
\]
Every positive integer has a unique factorization $n=rm$ with
$m\in\mathcal S(\mathcal P)$ and $(r,Q)=1$, where $Q=\prod_{p\in\mathcal P}p$.
Call $r=\rho_{\mathcal P}(n)$ its rough part. Choose distinct rough integers
$r_1,\ldots,r_L>N$ and real coefficients $v_{r_\ell}$. Extend them to
all positive indices by
\begin{equation}\label{eq:orbit-coeff}
 v_n=
 \begin{cases}
 v_{r_\ell},&\rho_{\mathcal P}(n)=r_\ell\text{ for some }\ell,\\
 0,&\text{otherwise}.
 \end{cases}
\end{equation}
Thus $v_n=0$ for $n\le N$ and $v_{pn}=v_n$ for every $p\in\mathcal P$.
Write
\begin{equation}\label{eq:orbit-source}
 K(u)=\sum_{n\ge1}v_n\phi_n(u),\qquad
 \psi_K(u)=\sum_{n\ge1}v_np_n(u),\qquad
 V(t,z)=\int_0^\infty e^{tu^2}K(u)\cos(zu)\dd u .
\end{equation}

Impose the given homogeneous measurement conditions on $(K,\psi_K)$,
and also impose $\psi_K'(0)=0$. These are finitely many homogeneous
linear equations in the $L$ coefficients. Taking $L$ larger than the
number of equations produces a nonzero solution. Normalize it so that
$\max_\ell|v_{r_\ell}|=1$. Then
\begin{equation}\label{eq:K-majorant}
 |K(u)|\le\sum_{n>N}\phi_n(u)\le\Phi(u)\qquad(u\ge0).
\end{equation}

We verify two necessary nondegeneracy facts. First, $K$ is not identically
zero. Otherwise $8K=\psi_K''-\psi_K=0$. The general solution is
$Ae^u+Be^{-u}$, and the superexponential decay as $u\to+\infty$ forces
$A=B=0$. If $n_0$ is the least index with $v_{n_0}\ne0$, divide
$e^{-u}\psi_K(u)=0$ by $e^{-\pi n_0^2e^{4u}}$ and let $u\to+\infty$.
The remaining terms tend to zero by dominated convergence, giving
$v_{n_0}=0$, a contradiction.

Second, some odd derivative of $K$ at zero is nonzero. If all vanished,
real analyticity would make $K$ even near zero and then, by analytic
continuation along the real line, on all of $\R$. Apply the consequence
of \eqref{eq:defect} with coefficients $v_0=0$, $v_{-n}=v_n$ and
$\psi_K'(0)=0$. Theorem~\ref{thm:rigidity} would force $v_n=0$, again a
contradiction. There is therefore a first nonzero odd endpoint derivative,
say $K^{(p-1)}(0)\ne0$ with $p$ even. Multiplication by the even function
$e^{tu^2}$ leaves this first odd derivative unchanged. At every fixed
time, Lemma~\ref{lem:endpoint} yields
\begin{equation}\label{eq:control-tail}
 V(t,x)=A_Kx^{-p}(1+O(x^{-2})),\quad
 (\log|V(t,x)|)_{xx}=\frac p{x^2}+O(x^{-4})>0
\end{equation}
for sufficiently large $x$, where $A_K=(-1)^{p/2}K^{(p-1)}(0)\ne0$.
The logarithmic derivative estimate follows by inserting all three
estimates \eqref{eq:endpoint-jets} into $V''/V-(V'/V)^2$.

\subsection{The additional input about high zeros}

The collision construction uses the following published arithmetic input,
which is separate from every packet estimate in this manuscript.

\begin{quote}
For each fixed $0<T\le1/2$, all sufficiently high zeros of
$H(T,\cdot)$ are real and simple, there are infinitely many such zeros,
and only finitely many zeros are nonreal.
\end{quote}

The high-zero assertions are consequences of
\cite[Theorem 1.5(i)--(iii) and its multiplicity count in \S9]{Polymath}.
Finiteness of the nonreal zeros also uses the global strip bound
$|\Im z|\le\sqrt{1-2T}$, obtained from the time-zero strip
$|\Im z|\le1$ by the de Bruijn theorem recorded in
\cite[Theorem 3.2]{Polymath}. The time-zero strip follows from
\eqref{eq:xi}, the Euler product in $\Re s>1$, and the reflection
$\xi(s)=\xi(1-s)$ evident in \eqref{eq:mellin-split}.
Thus the nonreal zeros lie in a compact rectangle and are finite by
isolatedness of zeros of a nonzero entire function. We use this published
input as a theorem; its arithmetic asymptotic proof is not reproduced.
No input about zeros at $T=0$ is assumed.

For clarity we derive the logarithmic-curvature consequence needed below.
The source estimate in Lemma~\ref{lem:regularity} implies, for $|z|\le R$,
\[
 |H(T,z)|\le C_T\int_0^\infty
      \exp\bigl((R+9)u-\tfrac{\pi}{2}e^{4u}\bigr)\dd u .
\]
Here $e^{Tu^2}$ was absorbed into half the superexponential factor,
with a constant depending on $T$. Split the remaining exponential decay
in half again. The integral of $\exp(-\pi e^{4u}/4)$ is finite, and the
maximum of $(R+9)u-\pi e^{4u}/4$ on $u\ge0$ is $O(R\log(R+2))$.
Consequently $\log\max_{|z|\le R}|H(T,z)|=O_T(R\log(R+2))$:
the entire function has order at most one.

We now apply the classical Hadamard factorization theorem in its standard
order-at-most-one form. If an entire function of this order is nonzero
at the origin, it is a product of its genus-one zero factors times
$e^{a+bz}$, and $\sum_\rho|\rho|^{-2}<\infty$. Since $H(T,0)>0$,
twice taking its logarithmic derivative gives, for real $x$ away from
zeros,
\begin{equation}\label{eq:logH-zeros}
 (\log|H(T,x)|)_{xx}=-\sum_\rho\frac1{(x-\rho)^2}.
\end{equation}
Multiplicity is included; the differentiated series converges locally
uniformly away from the zeros. The exponential factor has zero second
logarithmic derivative. Each real zero contributes a negative term.
A nonreal conjugate pair $\alpha\pm i\beta$ contributes
\[
 -\frac{2((x-\alpha)^2-\beta^2)}
            {((x-\alpha)^2+\beta^2)^2},
\]
which is negative once $x>\alpha+|\beta|$. There are only finitely many
such pairs. Hence
\begin{equation}\label{eq:logH-negative}
 (\log|H(T,x)|)_{xx}<0
\end{equation}
at every sufficiently large real $x$ that is not a zero.

\subsection{A finite-prime-compatible collision theorem}

\begin{theorem}[Scoped collision controls]\label{thm:controls}
Fix $0<T\le1/2$, a finite prime set $\mathcal P$, an integer $N$, and
finitely many real linear measurements as above. There exist bounded
strictly positive coefficient sequences $a_n^{(j)}$ and real numbers
$x_j\to+\infty$ such that:
\begin{enumerate}
\item $a_n^{(j)}=1$ for $n\le N$, and
      $a_{pn}^{(j)}=a_n^{(j)}$ for $p\in\mathcal P$ and all $n$;
\item the source and primitive have exactly the theta values for the
      specified measurements, including $\psi_{a^{(j)}}'(0)=-1/2$;
\item their heat transforms have a nondegenerate real double zero at
      $(T,x_j)$;
\item $\sup_n|a_n^{(j)}-1|\to0$, and their transforms and every fixed finite
      collection of spectral derivatives converge to $H$ uniformly on
      each set $0\le t\le T_+$, $|\Im z|\le Y$, with no restriction on
      $\Re z$.
\end{enumerate}
Every one of these sources fails the full reflection law
\eqref{eq:weighted-reflection} and fails evenness of its analytic source
on the whole real line.
\end{theorem}
\begin{proof}
Choose $K$ as in \eqref{eq:orbit-source}. For sufficiently large $x$,
$V(T,x)$ is nonzero. Let $(\beta_j,\beta_{j+1})$ be consecutive sufficiently
high real zeros of $H(T,\cdot)$, and define
$q(x)=H(T,x)/V(T,x)$ on the closed interval. This function is smooth,
vanishes at both endpoints, and has constant nonzero sign inside.
Equations \eqref{eq:control-tail} and \eqref{eq:logH-negative} imply
\[
 (\log|q(x)|)_{xx}<0 .
\]
The simple endpoint zeros give
$(\log|q|)_x\to+\infty$ at the left endpoint and
$(\log|q|)_x\to-\infty$ at the right endpoint. It is strictly decreasing,
so has a unique zero $x_j$ inside. At this point $q'(x_j)=0$ and
$q''(x_j)/q(x_j)<0$.

Set $\delta_j=-q(x_j)$, keep this scalar fixed as time varies, and put
\[
 a_n^{(j)}=1+\delta_jv_n,\qquad
 \mathcal H_j(t,z)=H(t,z)+\delta_jV(t,z).
\]
At time $T$, $\mathcal H_j=V(q+\delta_j)$. Therefore
\[
 \mathcal H_j(T,x_j)=\partial_x\mathcal H_j(T,x_j)=0,\qquad
 \partial_x^2\mathcal H_j(T,x_j)=V(T,x_j)q''(x_j)\ne0 .
\]
Both summands obey the heat equation, so this is a nondegenerate double
event of that equation. If desired, its local splitting follows by the
same Taylor calculation as in Section~\ref{sec:collisions}.

The rapid real-axis decay of $H(T,x)$, together with the power tail
\eqref{eq:control-tail}, gives $\delta_j\to0$ faster than any fixed
inverse power of $x_j$. Discard finitely many indices so that
$|\delta_j|<1$. Since $|v_n|\le1$, all $a_n^{(j)}$ are then positive.
Their finite prime identities, initial coefficients and measurements
follow directly from the construction of $v_n$ and its homogeneous
constraints. The perturbation is nonzero because $q(x_j)\ne0$.

For the convergence assertion, Lemma~\ref{lem:regularity} gives a finite
constant
\[
 M_{k,T_+,Y}=\int_0^\infty u^k e^{T_+u^2+Yu}
                         \sum_{n>N}\phi_n(u)\dd u .
\]
Differentiating the integral for $\mathcal H_j-H$ gives the uniform bound
$|\partial_z^k(\mathcal H_j-H)|\le|\delta_j|M_{k,T_+,Y}$ on the stated
set. Lemma~\ref{lem:strip-tail} makes this bound explicit whenever its
starting index is admissible; a finite initial segment can otherwise be
bounded directly. Mixed time derivatives follow as well, since each
time derivative inserts $u^2$.

Finally take $a_0^{(j)}=1$. If the full reflection law held,
Theorem~\ref{thm:rigidity} would force $a_n^{(j)}=1$ for every $n$,
contrary to the nonzero perturbation. If the analytic source were even,
the imposed primitive endpoint and \eqref{eq:defect} would imply that
same reflection law. Both assertions therefore fail.
\end{proof}

The meaning of the prime identities can also be checked at source level:
if $a_{pn}=a_n$, then
\[
 D_p\Phi_a=\sum_{n\ge1}a_n\phi_{pn}
          =\sum_{p\mid k}a_k\phi_k .
\]
Thus the finite selector algebra and the complete boundary formulas are
retained for every prime specified in the theorem.
The evenness of the cosine transform in $z$ is also retained. This is
different from evenness of the analytic arithmetic source in $u$.
Its even extension from $u\ge0$ has a nonmatching odd derivative at zero
and is not a smooth theta source on the whole line.

\subsection{Why the stronger packet bounds survive the controls}

Assume $2\in\mathcal P$ and write $\eta=|\delta_j|<1$.
Every positive coefficient of the controlled source lies between
$1-\eta$ and $1+\eta$. The positive-real-part theorem in
Section~\ref{sec:positivity} applies to these coefficients without a
perturbative loss. In particular, with
$\alpha^{(j)}=(\log A^{(j)})_x$ on the real axis,
\begin{equation}\label{eq:control-sharp-curvature}
 (\alpha^{(j)})^2+|\partial_x\alpha^{(j)}|\le\frac1{64}.
\end{equation}
If the remainder begins at $m$, its logarithmic bound is
$1/(4m^4)$ in the same form. Hence for $r^{(j)}=R^{(j)}/A^{(j)}$,
\[
 |r_x^{(j)}|\le d_m r^{(j)},\qquad
 |r_{xx}^{(j)}|\le d_m^2r^{(j)},\qquad
 d_m=\frac18+\frac1{2m^2}.
\]
These exact logarithmic constants therefore coexist with collisions.

Absolute relative-amplitude bounds require the coefficient comparison.
The controlled carrier has real part at least $(1-\eta)\Re H_2$ in
$|\Im z|\le1$, while the source, endpoint and moment upper bounds increase
by at most $1+\eta$. Put $k_\eta=(1+\eta)/(1-\eta)$. The real and
complex relative estimates of Section~\ref{sec:positivity} give
\[
 0<r^{(j)}(t,x)\le k_\eta\rho_m,\qquad
 |r^{(j)}(t,z)|\le k_\eta\rho_{s,m}\quad(|\Im z|\le1).
\]
Thus the absolute remainder-derivative tolerances acquire the same
$k_\eta$ factor; their original numerical values must not be asserted
for arbitrary positive coefficients. As $\eta\to0$ the amplitude
bounds approach the theta bounds. The controls also preserve any fixed
finite set of endpoint constraints included among the measurements.

The conclusion is confined to this coefficient-perturbation class.
The theorem prescribes the positive event time but neither a finite
event height nor an effective height threshold. It is not uniform as
$T\downarrow0$. It does not assign a Newman constant to the perturbed
families. In particular, it is not a construction of collisions for the
theta coefficients, and it does not exclude arguments that use the full
modular identity quantitatively.
